\documentclass[11pt,a4paper]{article}
\usepackage[utf8]{inputenc}
\usepackage[T1]{fontenc}
\usepackage[spanish]{babel}
\usepackage{geometry}
\usepackage{verbatim}
\geometry{margin=2.5cm}

\title{Analizador léxico y sintáctico de C-TDS \\ \large Diseño e interfaz del AST}
\author{Agustin Alieni, Francisco Natale y Julian Varea}
\date{}

\begin{document}
\maketitle

\section{Introducción}

En esta etapa construimos el analizador léxico (\texttt{lexer.l}, con Flex) y
el analizador sintáctico (\texttt{bison.y}, con Bison), y definimos la
interfaz del árbol de sintaxis abstracta (AST) para las etapas posteriores.
El AST todavía no está implementado: \texttt{ast.h} contiene su interfaz,
\texttt{ast.c} tiene los cuerpos pendientes y las acciones semánticas de
\texttt{bison.y} están comentadas.

\section{Analizador léxico}

El lexer reconoce las palabras reservadas (\texttt{boolean}, \texttt{int},
\texttt{float}, \texttt{void}, \texttt{if}, \texttt{else}, \texttt{while},
\texttt{return}), identificadores, constantes numéricas y booleanas,
operadores, delimitadores y comentarios.

\subsection{\texttt{main} no es palabra reservada}

\texttt{main} se reconoce como un identificador. La exigencia de que exista
un método \texttt{main} con una firma determinada es una regla semántica, por
lo que su verificación queda para esa etapa.

\subsection{Número de línea}

Activamos \texttt{\%option yylineno} y usamos el número de línea en los
errores léxicos y en \texttt{yyerror}.

\section{Decisiones de diseño}
\label{sec:decisiones}

\subsection{Declaraciones globales}

La primera versión separaba las declaraciones globales:

\begin{verbatim}
Program : VariableDeclarations MethodDeclarations ;
\end{verbatim}

Esto producía un conflicto shift/reduce porque las declaraciones de variables
y métodos comparten el prefijo \texttt{Type ID}. Por ejemplo:

\begin{verbatim}
int a;
int sumar(int a, int b) { }
\end{verbatim}

Al encontrar el segundo \texttt{int}, el parser no podía decidir todavía si
debía reducir una declaración de variable o continuar hacia un método.

Unificamos ambas listas:

\begin{verbatim}
Declarations
    : /* vacio */
    | Declaration Declarations
    ;

Declaration
    : VariableDeclaration
    | MethodDeclaration
    ;
\end{verbatim}

Así, variables y métodos pueden aparecer intercalados. También eliminamos el
no terminal \texttt{ReturnType} e integramos \texttt{Type} y \texttt{VOID}
directamente en \texttt{MethodDeclaration}.

Verificamos la gramática con \texttt{bison -d -v}: el reporte no presenta
conflictos shift/reduce ni reglas inalcanzables. También probamos el parser
con el ejemplo de la especificación y con variantes que intercalan variables
y métodos, incluyen condicionales, ciclos y llamadas a métodos.

\subsection{Precedencia de operadores}

Usamos \texttt{\%left} y \texttt{\%right} de Bison según la precedencia de la
especificación: disyunción, conjunción, igualdad, relacionales, suma y resta,
multiplicación/división/resto, y negación y menos unario.

El menos unario comparte el token \texttt{-} con la resta, pero tiene mayor
precedencia. Para distinguir ambos casos usamos un token ficticio:

\begin{verbatim}
%right NOT UMINUS
...
    | '-' Expression %prec UMINUS
\end{verbatim}

El lexer nunca devuelve \texttt{UMINUS}; solo se usa para asignar precedencia
a esa regla.

\subsection{Condición del \texttt{while}}

La especificación presenta \texttt{while} sin paréntesis, a diferencia de
\texttt{if}. Según lo acordado en clase, lo modelamos con la condición entre
paréntesis:

\begin{verbatim}
    | WHILE '(' Expression ')' Block
\end{verbatim}

\subsection{Comentarios multilínea}

Por el momento asumimos que los comentarios multilínea están correctamente
cerrados. El caso contrario queda como problema conocido del lexer.

\section{Diseño del AST}
\label{sec:ast}

\subsection{Interfaz}

Separamos el AST de \texttt{bison.y}: \texttt{ast.h} define sus tipos y
funciones constructoras, mientras que \texttt{ast.c} contiene sus
implementaciones.

En esta etapa \texttt{ast.c} mantiene los cuerpos vacíos porque todavía no
cerramos decisiones sobre manejo de memoria y sobre el uso de
\texttt{left}/\texttt{mid}/\texttt{right} y \texttt{children[]}.

Las acciones semánticas de \texttt{bison.y} están escritas pero comentadas
como referencia de cómo se construirá cada nodo.

\subsection{Estructura de \texttt{NodeAST}}

Cada \texttt{NodeAST} contiene:

\begin{itemize}
  \item \texttt{nodeType}: construcción sintáctica representada.
  \item \texttt{Symbol}: identificador o valor asociado, cuando corresponde.
  \item \texttt{left}, \texttt{mid}, \texttt{right}: hijos para estructuras
  con una cantidad fija y pequeña de elementos.
  \item \texttt{children[]}: hijos de cantidad variable.
  \item \texttt{type}: tipo semántico del nodo.
\end{itemize}

\texttt{type} se mantiene separado de \texttt{nodeType}: el primero representa
el tipo de dato y el segundo la construcción sintáctica.

\subsection{Criterio para crear nodos}

No creamos un nodo para cada no terminal. Solo se representan construcciones
que aportan información al AST. Por ejemplo, \texttt{Statement} agrupa
alternativas y simplemente propaga el nodo de la alternativa elegida.
\texttt{Block}, en cambio, tiene un nodo propio porque representa un scope y
agrupa declaraciones y sentencias.

\subsection{Listas de hijos}

\texttt{NodeList} se usa durante el parsing para acumular hijos en reglas
recursivas. \texttt{attachChildren} copia esa lista al arreglo
\texttt{children[]} del nodo padre. El AST final queda compuesto por
\texttt{NodeAST}, sin la lista auxiliar.

\section{Testing}
\label{sec:tests}

Agregamos tests unitarios para lexer y parser en \texttt{tests/}, usando
Unity. Se ejecutan con:

\begin{verbatim}
bash test_suite.sh
\end{verbatim}

Los tests del lexer cubren tokens mal reconocidos, literales numéricos
inválidos, operadores no permitidos, comentarios, símbolos inválidos y
números de línea.

Los tests del parser verifican programas válidos e inválidos, incluyendo
errores en punto y coma, llaves, firmas, sentencias, expresiones y
paréntesis. También cubren la intercalación de variables y métodos globales.

Los tests solo verifican reglas léxicas y sintácticas. No evalúan reglas
semánticas ni la precedencia de operadores, porque el AST todavía no está
implementado.

\section{Problemas conocidos}
\label{sec:problemas}

\subsection{Comentario multilínea sin cerrar}

La expresión regular actual para comentarios multilínea requiere que exista
el cierre \texttt{*/}. Si el archivo termina antes, el contenido se procesa
con las reglas siguientes y se reportan errores léxicos sobre sus caracteres.

Caso de test:

\begin{verbatim}
/* este comentario
nunca se cierra
int nuncaSeEjecuta;
\end{verbatim}

La solución prevista es usar un estado exclusivo de Flex, por ejemplo
\texttt{\%x COMMENT}, y una regla \texttt{<<EOF>>} que informe que el
comentario no fue cerrado. Queda pendiente de implementar.

\subsection{Alcance de esta entrega}

El AST queda limitado a su interfaz en \texttt{ast.h}. Las implementaciones de
\texttt{ast.c} y las acciones semánticas de \texttt{bison.y} quedan para la
siguiente etapa, una vez cerradas las decisiones de diseño pendientes.

\end{document}
